% -----------------------------------------------------------
\section{Hearken}
% -----------------------------------------------------------
	\label{HEARKEN}
	
% % ----------------------
% \subsection{See also}
% % ----------------------

% ----------------------
\subsection{1828 American Dictionary of the English Language} % *1828
% ----------------------
	\label{HEARKEN-1828}

	\begin{compactitem}
		\item[1] To listen; to lend the ear; to attend to what is uttered, with eagerness or curiosity.
		\item[2] To attend; to regard; to give heed to what is uttered; to observe or obey.
		\item[3] To listen; ; to attend; to grant or comply with.
	\end{compactitem}

% ----------------------
\subsection{Oxford English Dictionary} % *OED
% ----------------------
	\label{HEARKEN-OED}
	
	Etymology: ``Old English \textit{hercnian}, \textit{heorcnian}, \textit{hyrcnian}, formed with suffix \textit{-n-} < *\textit{heorci-an}, the Old English type of \textsc{hark} v.
	
	``The spelling harken, which agrees with that of \textsc{HARK} v., and is at once more regular and of earlier standing, is the accepted one in modern American Dictionaries, and is preferred by some good English writers; but in current English use it is much less frequent than hearken. The preference for the latter spelling is probably due to association with \textsc{HEAR} v., supported by the analogy of heart and hearth.''

	\begin{compactitem}
		\item[1] \textit{intransitive}. To apply the ears to hear; to listen, give ear.
		\item[2] \textit{intransitive}. To listen privily; to play the eavesdropper; to eavesdrop. \textit{Obsolete}.
		\item[3.a] \textit{intransitive}. To apply the mind to what is said; to attend, have regard; to listen with sympathy or docility. Const. \textit{to}.
		\item[3.b] with \textit{on}. \textit{Obsolete}.
		\item[4.a] transitive. To hear with attention, give ear to (a thing); to listen to; to have regard to, heed; to understand, learn by hearing; to hear, perceive by the ear. Now only \textit{poetic}.
		\item[4.b] With personal object (originally dative as in 1; but this afterwards levelled with the accusative or objective). \textit{Obsolete} exc. \textit{dialect}.
		\item[5] \textit{intransitive}. \textit{\textbf{hearken to}}: Listen, give ear.
		\item[6] \textit{intransitive}. To seek to hear tidings; to make inquiries, to inquire \textit{after}, ask \textit{for}. \textit{Obsolete}.
		\item[7] \textit{intransitive}. To lie in wait; to wait. \textit{Obsolete}.
		\item[8] transitive. To get to hear of; to search out or find by inquiry. \textit{Obsolete}.
		\item[9] \textit{intransitive}. To have regard or relation. \textit{Obsolete}. \textit{rare}.
		\item[10] To talk in one's ear, to whisper. \textit{Obsolete} exc. \textit{Scottish}.
	\end{compactitem}
	
% % ----------------------
% \subsection{Buck's Theological Dictionary (1833)} % *BTD
% % ----------------------
	% \label{HEARKEN-BTD}

	% \begin{compactdesc}
		% \item[PAGE\#]
	% \end{compactdesc}
	
% % ----------------------
% \subsection{Restoration Lexicon} % *RL *RESTORATION LEXICON
% % ----------------------
	% \label{HEARKEN-OED}

	% \begin{compactitem}
		% \item[]
	% \end{compactitem}

% % ----------------------
% \subsection{Scriptural Usage} % *SCRIPTURES
% % ----------------------
	% \label{HEARKEN-SCRIPTURES}

	% % -----
	% \subsubsection{Old Testament} % *OT
	% % -----
		% \label{HEARKEN-OT}
	
		% % --
		% \paragraph{KJV}
		% % --
			% \label{HEARKEN-OT-KJV}
	
			% \begin{tabular}{ll}
				% Total occurrences in the OT & 
			% \end{tabular}
			
			% \begin{compactdesc}
				% \item[]
			% \end{compactdesc}
			
		% % --
		% \paragraph{Hebrew Bible}
		% % --
			% \label{HEARKEN-HEBREW}
		
			% %
			% \subparagraph{\texorpdfstring{\he{sample word}}{}} % paste actual hebrew text here
			% %
				% \label{PAGA} % whatever is in step bible without periods
			
				% \begin{compactdesc}
					% \item[HALOT , PDF ] \textbf{} % Use the 1994 PDF
						% \begin{compactitem}
							% \item[1]
						% \end{compactitem}
					% \item[BDB , PDF ] \textbf{}
						% \begin{compactitem}
							% \item[1]
						% \end{compactitem}
					% \item[DCH PDF ] \textbf{} % don't include the actual page number, they repeat from volume to volume
						% \begin{compactitem}
							% \item[1]
						% \end{compactitem}
				% \end{compactdesc}
				
				% \begin{compactdesc}
					% \item[Some OT uses] \textbf{}
						% \begin{compactdesc}
							% \item[]
						% \end{compactdesc}
				% \end{compactdesc}
			
		% % --
		% \paragraph{Septuagint}
		% % --
			% \label{HEARKEN-LXX}
		
			% %
			% \subparagraph{\texorpdfstring{\gr{sample word}}{}}
			% %
		
	% % -----
	% \subsubsection{New Testament} % *NT
	% % -----
		% \label{HEARKEN-NT}
	
		% % --
		% \paragraph{KJV}
		% % --
			% \label{HEARKEN-NT-KJV}		
	
			% \begin{tabular}{ll}
				% Total occurrences in the NT & 
			% \end{tabular}
			
			% \begin{compactdesc}
				% \item[]
			% \end{compactdesc}
			
		% % --
		% \paragraph{Greek New Testament} % *GREEK
		% % --
			% \label{HEARKEN-GREEK}
		
			% %
			% \subparagraph{\texorpdfstring{\gr{sample word}}{}}
			% %
				% \label{KATALLAGE}
			
				% \begin{compactdesc}
					% \item[BDAG ] \textbf{}
						% \begin{compactitem}
							% \item 
						% \end{compactitem}
					% \item[CGL , PDF ] \textbf{}
						% \begin{compactitem}
							% \item 
						% \end{compactitem}
					% \item[LSJ , PDF ] \textbf{}
						% \begin{compactitem}
							% \item 
						% \end{compactitem}
					% \item[Brill , PDF ] \textbf{}
						% \begin{compactitem}
							% \item 
						% \end{compactitem}
				% \end{compactdesc}
				
				% \begin{tabular}{ll}
					% Total occurrences in the NT & 
				% \end{tabular}
				
				% \begin{compactdesc}
					% \item[Some NT uses] \textbf{}
						% \begin{compactdesc}
							% \item[]
						% \end{compactdesc}
				% \end{compactdesc}
				
			% %
			% \subparagraph{TDNT: \texorpdfstring{\gr{sample word}}{}  (3:536--556)}
			% %
				% \label{KATALLAGE-TDNT}
		
	% % -----
	% \subsubsection{Book of Mormon} % *BOM
	% % -----
		% \label{HEARKEN-BOM}
	
		% \begin{tabular}{ll}
			% Total occurrences in the BOM & 
		% \end{tabular}
		
		% \begin{compactdesc}
			% \item[]
		% \end{compactdesc}
		
	% % -----
	% \subsubsection{Doctrine and Covenants} % *DC
	% % -----
		% \label{HEARKEN-DC}
	
		% \begin{tabular}{ll}
			% Total occurrences in the DC & 
		% \end{tabular}
		
		% \begin{compactdesc}
			% \item[]
		% \end{compactdesc}
		
	% % -----
	% \subsubsection{Pearl of Great Price} % *PGP
	% % -----
		% \label{HEARKEN-PGP}
	
		% \begin{tabular}{ll}
			% Total occurrences in the PGP & 
		% \end{tabular}
		
		% \begin{compactdesc}
			% \item[]
		% \end{compactdesc}
		
% % ----------------------
% \subsection{Key Phrases} % *KEYPHRASES *PHRASES
% % ----------------------
	% \label{HEARKEN-PHRASES}

	% \begin{compactdesc}
		% \item[] \textbf{\#times} | list
			% % \begin{compactdesc}
				% % \item[Bible anchor verses] \phantom{.}
					% % \begin{compactdesc}
						% % \item[Romans 5:5] \gr{}
					% % \end{compactdesc}
				% % \item[Commentary] ---
			% % \end{compactdesc}
	% \end{compactdesc}
	
% % ----------------------
% \subsection{Bible Dictionaries} % *DICTIONARIES
% % ----------------------
	% \label{HEARKEN-BD}

% % ----------------------
% \subsection{Restoration Usage} % *RESTORATION
% % ----------------------
	% \label{HEARKEN-RESTORATION}

	% % -----
	% \subsubsection{Joseph Smith} % *JS
	% % -----
		% \label{HEARKEN-JS}
	
		% \begin{compactdesc}
			% \item[{\hyperref[KEY]{Date/Source}}]
		% \end{compactdesc}
		
	% % -----
	% \subsubsection{Temple Ordinances}
	% % -----
		% \label{HEARKEN-TEMPLE}
	
		% \begin{compactdesc}
			% \item[{\hyperref[KEY]{Date/Source}}]
		% \end{compactdesc}
	
	% % -----
	% \subsubsection{Brigham Young} % *BY
	% % -----
		% \label{HEARKEN-BY}
	
		% \begin{compactdesc}
			% \item[{\hyperref[KEY]{Date/Source}}]
		% \end{compactdesc}
	
% % ----------------------
% \subsection{Commentary and Studies} % *COMMENTARY *STUDIES
% % ----------------------
	% \label{HEARKEN-COMMENTARY}

	% % -----
	% \subsubsection{Key Points}
	% % -----
		% \label{HEARKEN-KEYS}
	
		% \begin{compactitem}
			% \item
		% \end{compactitem}
		
	% % -----
	% \subsubsection{Sample Study}
	% % -----
		% \label{HEARKEN-study}